\documentclass[conference]{IEEEtran}
\IEEEoverridecommandlockouts
\usepackage{cite}
\usepackage{amsmath,amssymb,amsfonts}
\usepackage{algorithmic}
\usepackage{graphicx}
\usepackage{textcomp}
\usepackage{xcolor}
\usepackage[ruled,vlined]{algorithm2e} 
\usepackage{url}
\usepackage[utf8]{inputenc}
\usepackage{textgreek}
\usepackage{array}
\usepackage{listings}
\usepackage{float}
\usepackage{booktabs}
\usepackage{fancyhdr}
\usepackage{tikz}
\usepackage{hyperref}
\usepackage{comment}

\usetikzlibrary{arrows.meta, positioning}
\usepackage{xurl}
\hypersetup{
    breaklinks=true}

\def\BibTeX{{\rm B\kern-.05em{\sc i\kern-.025em b}\kern-.08em
    T\kern-.1667em\lower.7ex\hbox{E}\kern-.125emX}}

% % === Custom Page Styles for ICoACT 2025 ===
% \fancypagestyle{firstpage}{
%   \fancyhf{} % Clear header and footer
%   \fancyfoot[L]{%
%     \begin{minipage}[t]{\textwidth}
%       \normalsize{\textcolor{black}{979-8-3315-6719-4/25/\$31.00 \textcopyright{} 2025 IEEE}}
%     \end{minipage}%
%   }
%   \renewcommand{\headrulewidth}{0pt}
%   \renewcommand{\footrulewidth}{0pt}
% }

% \pagestyle{fancy}
% \fancyhf{}
% \renewcommand{\headrulewidth}{0pt}
% \renewcommand{\footrulewidth}{0pt}
    
\begin{document}

\title{A Multimodal Retrieval-Augmented Framework for Real-Time Disaster Intelligence with Temporal Grounding and Safety Validation}

\author{
\IEEEauthorblockN{
Christopher Aseer J Albert\textsuperscript{1},
Ashutosh Monish Bhonsle\textsuperscript{2},
Prisha Andrew Bhosale\textsuperscript{3},
Harshvardhan Bhota\textsuperscript{4},\\
Chaitanya Vijaykumar Mahamuni\textsuperscript{5}
}

\IEEEauthorblockA{
\textsuperscript{1,2,3,4}\textit{Department of Computer Engineering, Fr. C Rodrigues Institute of Technology, Vashi, India}\\
\textsuperscript{5}\textit{Electronics and Telecommunication Engineering, University of Mumbai, Mumbai, India}\\
\textsuperscript{1}\texttt{christpheraseer5080@gmail.com}, \textsuperscript{2}\texttt{bhonsle.ashutosh@gmail.com}, 
\textsuperscript{3}\texttt{prisha.bhosale22@gmail.com},\\
\textsuperscript{4}\texttt{harshvardhanbohta@gmail.com}, \textsuperscript{5}\texttt{chaitanyamahamuni91@gmail.com}
}
}

\maketitle
% \thispagestyle{firstpage}

\begin{abstract}


Rapid and accurate situational awareness is critical for effective disaster response and emergency decision-making. Modern disaster scenarios generate large volumes of heterogeneous data, including satellite and unmanned aerial vehicle (UAV) imagery, crisis-related textual information, environmental sensor readings, and official disaster management documents. However, existing multimodal disaster AI systems process modalities in isolation, lack temporal grounding between historical and live data, and do not enforce deterministic safety checks on generated responses. To address this limitation, this paper proposes a unified multimodal Retrieval-Augmented Generation (RAG) framework for disaster response.  The proposed system integrates visual, textual, sensor, and domain knowledge through modality-specific embedding and retrieval, followed by evidence fusion and knowledge-grounded generation. A safety-aware verification mechanism is incorporated to mitigate hallucinations and improve the reliability of generated outputs. The framework is evaluated using a combination of real-world and simulation-based disaster data, demonstrating its effectiveness in enhancing situational awareness and providing explainable, grounded disaster intelligence. We introduce (i) a Temporal Grounding Injection mechanism to separate historical and live evidence, and (ii) a Deterministic Safety Enforcement Layer that guarantees operator-safe outputs independent of model behavior.
\end{abstract}

\begin{IEEEkeywords}
Disaster Response; Multimodal Learning; Retrieval-Augmented Generation; Situational Awareness; Large Language Models; Safety-Critical AI
\end{IEEEkeywords}

\section{Introduction}

Earthquakes, floods, landslides, and hurricanes are natural phenomena that lead to massive amount of death of human beings, destruction of infrastructure and disturbance of the socio-economic life across the globe ~\cite{b1}. The rapid and precise situational awareness is a key to the successful disaster response to assist in emergency decision making and search-and-rescue activities ~\cite{b2}. The contemporary disaster environment produces massive amounts of non-standardized information on sources like satellite and unmanned aerial vehicle (UAV) imagery, social media, and environmental sensors, and official disaster management policies. Nevertheless, the problem of successfully synthesizing and reasoning with these disparate multimodal data streams is still a significant challenge with regard to the current disaster response systems ~\cite{b3}.

The latest advances in the field of artificial intelligence (AI) made it possible to better monitor disasters by evaluating the damage based on the satellite and UAV images and analyzing social media related to the crisis and emergency messages with the help of natural language processing ~\cite{b4}. The sensor driven methods also play a role by identifying the seismic, hydrological and environmental variations in real time. Although the multimodal learning techniques have been studied to have a synthesis of visual and textual information, most of the current techniques are based on the classification or detection processes and lack of contextual reasoning, explainability and incorporation of external domain knowledge, which is also crucial in a disaster recovery situation that deals with safety aspects.

Large Language Models (LLMs) are strong reasoners with generating skills, but can be affected by hallucinations when applied without grounding mechanisms, especially in the high-risk domain ~\cite{b5}. Retrieval-Augmented Generation (RAG) is a system that provides answers to this limitation by basing responses on external knowledge by means of retrieval; but existing systems of RAG are highly text-based and are rarely used in multimodal disaster intelligence ~\cite{b6}. In addition, simulated and synthetic data is commonly used due to limited or unavailability of real-world data of disasters, which are challenging to analyze under controlled conditions. Such restrictions emphasize the necessity of such a cohesive and multimodal retrieval-augmented system that can combine visual, textual, sensor, and domain information to ensure dependable and explicable disaster response.

To overcome the above challenges, the paper presents a single multimodal Retrieval-Augmented Generation framework of disaster response and proposes how Crisis-focused textual data, satellite and UAV image, sensor-based signals, and authoritative disaster management documents can be incorporated into a single retrieval-driven reasoning pipeline. The proposed system conducts modality embedded and indexing, and then multimodal recovery, evidence fusion, and knowledge-informed generation to provide efficient and accurate, explainable, and safety-aware disaster intelligence.

The primary contributions of this work are as follows:

\begin{enumerate}
    \item An integrated multimodal Retrieval-Augmented Generation system of disaster response that incorporates visual, textual, sensor and domain knowledge in a single reasoning pipeline.
    
    \item A retrieval and fusion system that is modality aware and allows contextual reasoning of the heterogeneous disaster data sources.
    
    \item A safety conscious generation and verification pipeline that would reduce the impact of hallucinations and provide trustworthy outputs, in high stakes disaster context.
\end{enumerate}



The remainder of this paper is organized as follows: Section II reviews related work, Section III details the methodology, Section IV describes implementation, Section V presents results and analysis, and Section VI concludes with future research directions.


\section{Literature Survey}

\subsection{Visual and Multimodal Situational Awareness for Disaster Response}
Satellite and UAV imagery damage assessment is a very important disaster management capacity. Although previously, the main method was based on convolutional neural networks, recent research has proven the usefulness of Vision Transformer architectures to high-resolution remote sensing images.~\cite{b15} As an example, Singh and Hoskere (2023) have indicated that transformer-based models can reach the state-of-the-art results on ultra-high-resolution aerial datasets in post-disaster damage classification. Nevertheless, the majority of available methods use visual data only or use weakly coupled multimodal pipelines, and they are not able to utilize contextual ground reports like social media information or field observations.



Most past systems of multimodal disaster response used late-fusion architectures, where textual and visual modalities are represented and reasoned about separately, which limits deep cross-modal representation learning and cross-modal reasoning ~\cite{b16}. In the recent past, unified multimodal foundation models have been introduced to overcome these shortcomings. An example of this paradigm is Emu3 ~\cite{b17}, which trains one transformer to reason multimodally in both spatiotemporal sequences, by predicting interleaved vision and language tokens. These unified models propose the possibility of joint situational insights through joint modelling of dynamics of visual scenes and semantic text information, transcending the modularity of traditional pipeline models.

\subsection{Integrating Social Media and Sensor Data for Disaster Response}

In addition to physical damage, social media text offers an essential source of real-time situational intelligence when there is a disaster, even based on satellite and UAV images. Nevertheless, this type of data is usually messy and unorganized. Researchers use Natural Language Processing (NLP) models to delete irrelevant information, as well as to categorize humanitarian needs, to draw a useful insight. As an example, Alam et al. (2021) presented the HumAID dataset and compared deep learning models to classify disaster-related tweets and proved the possibility of extracting actionable information by analyzing social media streams.~\cite{b18} 


The biggest problem in multimodal social media analysis is discrepancy between visual and textual communication, where images and text can communicate incompatible or loosely associated information. Abavisani et al. (2020) reviewed the multimodal fusion methods of categorizing crisis events and noted that learning about cross models representation is strong when addressing the issue at hand.~\cite{b19} Moreover, recent surveys highlight the increased importance of IoT-based sensing infrastructures in disaster management to monitor and take decisions in real-time. ~\cite{b20} A combination of social media analytics and physical sensor information can enhance the credibility and validation of situational reports and decrease information saturation among emergency responders.

\begin{comment}
\subsection{Integrating Textual and Visual Cues for Disaster Event Analysis}

Multimodal analysis of social media has become an increasingly effective approach for real-time disaster event classification and prioritization, as it allows textual descriptions and visual evidence shared during crises to be considered together. Recent frameworks combine transformer-based language models with deep vision networks to support tasks such as disaster detection, event type classification, semantic intent analysis, and visual damage assessment, achieving strong performance across multiple benchmarks. In particular, late-fusion architectures that integrate parallel text and image pipelines using weighted decision strategies have shown clear improvements over unimodal approaches, enabling more accurate prioritization of critical events for emergency response. However, these systems commonly operate under fixed, closed-world assumptions and fuse modalities only at the decision level, which limits deeper cross-modal reasoning and reduces their ability to adapt to unseen or evolving disaster scenarios. As a result, despite advances in automated disaster monitoring, their robustness in dynamic real-world settings remains limited.~\cite{b6}

\subsection{Disaster Event Identification Using Multimodal Data}

Situational awareness in disaster response increasingly draws on social media content, where textual descriptions and visual evidence together provide valuable insights into unfolding events. Previous research has shown that multimodal deep learning approaches can improve disaster information classification by effectively integrating complementary cues from text and images. However, many existing methods rely on a closed-world assumption, presuming that all disaster categories encountered during deployment are present in the training data. This assumption limits their practical applicability, as real-world disaster scenarios often involve unforeseen or evolving events. To overcome this limitation, recent studies have explored open-world learning frameworks that enable models to identify and appropriately handle unseen classes. In multimodal contexts, such approaches combine cross-modal feature fusion, attention-based reasoning, and unknown class detection mechanisms to enhance robustness. While these methods significantly improve decision-making in dynamic environments, further research is needed to ensure consistent reliability under rapidly changing disaster conditions.~\cite{b5}
\end{comment}

\subsection{Knowledge Graph and Retrieval-Augmented Generation Approaches for Disaster Events}
Knowledge-based decision support plays an important role in earthquake emergency management, where timely access to structured and reliable information is critical. Recent work has shown that knowledge graphs can effectively organize complex disaster-related knowledge, while large language models enable flexible information retrieval and reasoning. Building on these ideas, this study presents a framework that combines ontology-driven knowledge graph construction with retrieval-augmented generation to support earthquake emergency response. The approach integrates information from authoritative documents to build a coherent knowledge graph and uses hybrid retrieval strategies to improve the accuracy and contextual relevance of generated responses. While the framework demonstrates strong potential for text-based decision support, its reliance on textual sources limits its ability to incorporate multimodal data and adapt to rapidly evolving disaster situations. ~\cite{b7}

\begin{comment}
\subsection{Disaster Situational Awareness from Heterogeneous Data}

Situational awareness during disasters needs to incorporate heterogeneous data including sensor networks, geospatial imagery and ground-sourced reports to facilitate timely decision making. Previous literature demonstrates the significance of situational frameworks that select and control various data streams to minimize information overload and enhance its applicability to responders. Nonetheless, current methods typically handle the separate modalities and use late-stage aggregation, but not multi-modal reasoning. This constrains the utilization of complementary cues and the formulation of thorough and sound situational awareness in intricate disaster conditions.~\cite{b9}
\end{comment}

\subsection{Reliability, Safety, and Evaluation in Disaster AI}

Situational awareness during disasters needs to incorporate heterogeneous data including sensor networks, geospatial imagery and ground-sourced reports to facilitate timely decision making. Previous literature demonstrates the significance of situational frameworks that select and control various data streams to minimize information overload and enhance its applicability to responders. Nonetheless, current methods typically handle the separate modalities and use late-stage aggregation, but not multi-modal reasoning. This constrains the utilization of complementary cues and the formulation of thorough and sound situational awareness in intricate disaster conditions.~\cite{b10}


\subsection{Hallucinations in LLMs Visual Spatial Descriptions}
Hallucinations continue to be a major issue to multimodal large language models in visual spatial description tasks, in which the models tend to produce erroneous object relations or describe non-existent objects. In previous studies, it has been shown that these fallacies are specifically common in cases of ambiguous or partial visual input, which results in inaccurate spatial reasoning. The solutions to this are retrieval-augmented and hallucination-aware correction mechanisms, which are shown to be more accurate in a spatial description. These results emphasize that vision-language systems require clear mitigation of hallucinations. ~\cite{b14}
% \section{Methodology}


% \section{Implementation}
% \label{sec:implementation}


\section{Methodology}

This research is based on a modular multimodal Retrieval-Augmented Generation (RAG) framework. The Augmented Generation (RAG) and agentic reasoning architecture. The system is specially designed for complex The data relating to disaster were collected and the dependency on heavy clouds was reduced. Avoiding external data transfers; ensuring data is processed locally whenever feasible, all for data privacy and cost benefits. efficiency.

\subsection{Framework Overview: System Architecture}

The system's essence is controlled through a ReAct system. (Reason + Act) loop using LangGraph, powered by the Groq Llama-3.3-70B-versatile Large Language Model (LLM). To engage in multi-turn, context-aware interactions, conversation history is preserved across sessions utilizing LangGraph’s MemorySaver checkpointer.

\begin{figure}[htbp]
\centering
\includegraphics[width=\columnwidth]{figure1.jpeg}
\caption{Proposed Multimodal Agentic RAG Architecture. The framework ingests heterogeneous data streams, processes them locally for privacy, and orchestrates retrieval and response generation through a LangGraph ReAct agentic loop.}
\label{fig:architecture}
\end{figure}

\subsubsection{Hardware Specifications and Local Processing} 
A core, This leaves data as the methodological priority of this framework. privacy, removing cloud cost for the local processing modules.

\begin{itemize}
    \item \textbf{Using a vector database and embeddings:} Vector Database and Embeddings are the backbone of the system. the HuggingFace \texttt{BAAI/bge-small-en-v1.5} embedding model. Generating 384-dimensional embeddings, the selected model was specifically chosen for its efficiency; runs on a standard CPU, processes batches in can do $\sim$100ms without using GPU acceleration.
    \item \textbf{Utilization of the pipeline for vision to text (vision AI processing):} Runs a local instance of Ollama as the LLaVA 7B model. It only requires one-time download of 4.7GB. model can be run flexibly on the available local GPU or CPU. resources. That is to say, this architecture makes sure that sensitivity to the disaster is considered. imagery is processed locally only, without sharing with third-party. Party cloud providers will have to adopt a strict data governance approach.
\end{itemize}

\subsubsection{Mathematical Formulation of Semantic Retrieval}
To perform localised knowledge retrieval without cloud overhead, the system runs with \texttt{BAAI/bge-small-en-v1.5} model vector space.map documents and user queries to a 384 dimensional vector space. vector space. The similarity between a user query $q$ and a The cosine is computed on document $d$ in the corpus $D$. similarity. Due to the embedding nature of the model, the following notes apply: Compute L2-normalized vectors and retrieve them. Short-circuits to an extremely effective point-wise product operation:

\begin{equation}
S(q, d) = \sum_{i=1}^{384} E(q)_i \cdot E(d)_i
\end{equation}

Here $E(\cdot)$ denotes the embedded vector generated. This normalized configuration reduces the CPU cycle cost of during the retrieval phase, for fast inference on local hardware.

\subsection{Multimodal Data Ingestion and Volumetrics}
To create a coherent picture to be used in an operational context, The ingestion pipeline ingests four modalities of data: A single vector storage database. The system helps to target a heavily filtered SNR-optimized, subset of multimodal data over sheer volume. In the case of sensors with live thresholds, the number of total embedded documents can vary between approximately 3,195 to 5,162.

\begin{table}[htbp]
\centering
\caption{Ingestion Data Sources and Volumetrics}
\label{tab:volumetrics}
\renewcommand{\arraystretch}{1.3}
\begin{tabular}{@{}p{1.2cm}p{2.5cm}p{3.5cm}@{}}
\toprule
\textbf{Modality} & \textbf{Data Source} & \textbf{Embedded Volume} \\ \midrule
\textbf{Text} & CrisisLexT26 (2013) & 847 -- 1,500 docs \\
\textbf{Imagery} & xView2 Dataset & 50 structured docs \\
\textbf{Knowledge} & FEMA \& CDC Manuals & 3,000 -- 4,200 chunks \\
\textbf{Telemetry} & Live USGS IoT Sensors & $\sim$30 -- 40 (Volatile) \\ \bottomrule
\end{tabular}
\end{table}

\textbf{Meta Data Schema:} To help the agent to be able to accurately Remember to filter and weigh context for each of the documents embedded. ChromaDB follows a single metadata contract (JSON) A schema with modality, source, timestamp and location). This tagging is essential for downstream temporal. and spatial reasoning.

\subsubsection{The Knowledge base ingestion process and Hyperparameter Justification}
cation: For the official protocol data, text is processed using a hierarchical recursive character text splitter. The chunking The size of $C_{size} = 1000$ was explicitly set hyperparameters. characters that have a $C_{overlap} = 200$ characters. This configuration optimizes semantic coherence without resorting to the use of Truncation in the maximum context window of 512 tokens. the bge-small model. The 20\% overlap prevents critical loading rates from exceeding the limits. The procedural steps from being broken down over borders.

\begin{algorithm}[htbp]
\caption{Agentic API Execution \& Fault Tolerance Loop}
\label{alg:fallback}
\SetAlgoLined
\KwIn{User Query $Q$, Target Location $L$}
\KwOut{Formatted Operational Response $R$}
 Initialize ReAct state with query $Q$\;
 \eIf{$Q$ requires live telemetry}{
  Attempt API\_Call($L$)\;
  \uIf{API\_Call returns HTTP\_Success}{
   $Context \leftarrow$ Parse\_Data(API\_Call)\;
  }
  \uElseIf{Error == "Location Not Found"}{
   Log "Unrecognized geocode for $L$"\;
   $Context \leftarrow$ API\_Call(Fallback\_Location)\;
   Notify\_User(Fallback\_Triggered)\;
  }
  \ElseIf{Error == "API Unreachable"}{
   $Context \leftarrow$ Vector\_Search($Q$, KB\_Historical)\;
   Notify\_User(Historical\_Fallback)\;
  }
 }{
  $Context \leftarrow$ Vector\_Search($Q$, DB)\;
 }
 $R \leftarrow$ LLM\_Synthesize($Context$, $Q$)\;
 \Return{$R$}\;
\end{algorithm}

\subsection{The agentic reasoning and orchestration of tools task.}
The task of agentic reasoning and orchestration of tools. This system requires an active RAG system, unlike the passive RAG systems. An autonomous agent able to dynamically orchestrate specialized tools.

\begin{table}[htbp]
\centering
\caption{Agent Tool Repertoire}
\label{tab:tools}
\renewcommand{\arraystretch}{1.3}
\begin{tabular}{@{}ll@{}}
\toprule
\textbf{Tool Name} & \textbf{Functionality} \\ \midrule
\texttt{fetch\_live\_weather} & Retrieves real-time weather data. \\
\texttt{query\_water\_sensor} & Accesses live USGS river gauge heights. \\
\texttt{search\_official\_protocols} & Semantic search of FEMA/CDC docs. \\
\texttt{search\_social\_reports} & Search of local KB datasets. \\
\texttt{analyze\_disaster\_image} & On-the-fly LLaVA visual processing. \\ \bottomrule
\end{tabular}
\end{table}

\subsubsection{Large Language Model Configuration}
The cognitive This is a llama-3.3-70b-versatile powered engine. model. To be able to call the tools deterministically, the temperature $T = 0.1$ is strictly minimized (hyperparameter). Other parameters remain unchanged to prevent overfitting. the distribution of the output to certain scenarios.

\subsection{The Temporal Grounding Injection}
One of the major architectural problems is the solution of temporal grounding injection (TGI) Discrepancy in time with embedded historical data against live data. operational feeds. Instead of using complicated metadata, try to keep it simple. We define this mechanism as Temporal Grounding Injection (TGI), a 
retrieval-time transformation that injects explicit temporal labels into document context before LLM reasoning. We propose the concept of “Temporal Grounding Injection” to the LLM's reasoning space. via prompt engineering.

The search tool parses the during retrieval phase. automatically generate the metadata that describes the timestamp for each vector result $d$ and dynamically Injects a string-based temporal tag $T_{tag}$ before a string Add input to the LLM's context window:

\begin{equation}
C(d) = T_{tag}(t_d) \oplus d
\end{equation}

If $T_{tag}$ is equal to \texttt{[HISTORICAL — YYYY]} for Ths is a set of older corpora data and \texttt{[LIVE/RECENT]} for active sensor telemetry. The agent focuses on active for live telemetry. Studies of historical patterns in decision-making while explicitly projecting historical heuristics Assesses current scenarios with tics.

% \subsection{Generating and Using Agents}
% To promote robustness in disaster situations where remote APIs can break; agent is designed to deal with graceful degradation. The tool catches errors returned by external API. the exception, and returns a formalized text string to the agent, This is followed by an automatic fallback to a hardcoded sequence.



% ============================================================
% IEEE SECTION V: RESULTS AND DISCUSSION
% Paste this block into your IEEE paper.
% Requires: \usepackage{booktabs,amsmath,graphicx,multirow}
% ============================================================

% ============================================================
% SECTION V: RESULTS AND DISCUSSION
% IEEE two-column format
% Packages required: booktabs, amsmath, graphicx
% ============================================================
\section{IMPLEMENTATION}

\subsection{Datasets Used}

The experimental evaluation integrates four heterogeneous data modalities, each sourced
from real-world disaster datasets or live operational APIs. No synthetic or mock data
was used during ingestion or evaluation; the pipeline explicitly rejects fallback or
placeholder content to preserve data authenticity.

1) Social Media Text Corpus --- CrisisLex T26.

The CrisisLex T26 dataset~\cite{b21} was used as the primary source of disaster-related
social media data. The dataset contains crisis-oriented tweets collected from 26 real
disaster events. For this study, the 2013 Colorado Floods subset was selected, and
up to 500 tweets labelled ``Related and informative'' were retained after applying
the dataset relevance filters. Irrelevant and duplicate tweets were discarded prior
to indexing.

To improve temporal reasoning, tweet creation timestamps were reconstructed directly
from the Twitter Snowflake identifiers by shifting the 64-bit integer representation
22 bits to the right and adding the Twitter epoch offset of 1,288,834,974,657 ms.
This produced precise UTC timestamps unavailable in the raw dataset metadata. All
processed tweet documents were indexed under the \texttt{text} modality within the
vector database.

2) Satellite and Aerial Imagery Descriptions --- xView2.

Post-disaster visual intelligence was derived from the xView2 dataset~\cite{b22},
which contains high-resolution satellite imagery spanning floods, hurricanes,
wildfires, earthquakes, and volcanic disasters. Since raw imagery cannot be directly
embedded into a semantic vector store, a Vision-Language Model (VLM) pipeline was
implemented using Ollama LLaVA 7B running locally.

Fifty images were streamed directly from the compressed xView2 archive and processed
individually without full extraction to disk. Each image was analysed using a
structured prompting template that enforced four output components: SUMMARY,
DAMAGE, ACCESS, and FULL\_DESC. This approach ensured consistent machine-readable
descriptions suitable for downstream retrieval operations.

The generated descriptions were indexed under the
\texttt{imagery\_description} modality. In the event of Ollama unavailability, the
pipeline terminates rather than generating synthetic fallback descriptions, ensuring
that all visual embeddings remain grounded in authentic observations.

3) IoT Sensor Data --- USGS National Water Information System.

Hydrological telemetry data was retrieved from the United States Geological Survey
(USGS) National Water Information System (NWIS) Instantaneous Values API using the
endpoint \texttt{waterservices.usgs.gov/nwis/iv/}. The primary monitoring site was
gauge 06730200 (Boulder Creek at North 75th Street, Boulder, Colorado), selected due
to its historical association with the 2013 Colorado Floods.

Seven days of 15-minute interval gage-height readings corresponding to parameter code
\texttt{00065} were downloaded and cached locally in CSV format during database
construction. The ingestion pipeline embedded the 24 most recent measurements to
represent live conditions, while the top-10 WARNING and CRITICAL flood readings from
the seven-day interval were additionally indexed to preserve historical flood trends.

Flood severity labels were assigned automatically using the 10.0 ft flood-stage
threshold defined for the monitoring site. Readings were categorized as NORMAL,
WARNING, DANGER, or CRITICAL according to threshold exceedance percentage. All
sensor records were indexed under the \texttt{telemetry} modality.

4) Official Emergency Protocols --- FEMA/CDC Knowledge Base.

Three official emergency-management manuals were integrated into the knowledge base:
the CDC Preparedness and Response Capabilities manual, the FEMA Disaster Preparedness
Training Manual, and the FEMA Field Operations Guide.

PDF parsing was performed using PyPDF, after which documents were segmented using
\texttt{RecursiveCharacterTextSplitter}. Chunking was configured with a chunk size
of 1,000 characters and an overlap of 200 characters using a hierarchical separator
structure consisting of paragraph, newline, sentence, and whitespace boundaries.
This strategy preserved semantic continuity while remaining within the embedding
model’s token constraints.

All extracted chunks were indexed under the \texttt{knowledge\_base} modality,
yielding approximately 3,000--5,000 semantic fragments. A modality-filtering policy
was additionally enforced to prevent operational FEMA or CDC guidance from being
confused with social-media observations during retrieval.

Table~\ref{tab:kb_comp} summarizes the composition of the final vector knowledge base.

\begin{table}[h]
\centering
\caption{Knowledge Base Composition}
\label{tab:kb_comp}

\resizebox{\columnwidth}{!}{
\begin{tabular}{|l|l|l|l|}
\hline
Modality Tag & Source & Documents & Notes \\ \hline

\texttt{text} &
CrisisLex T26 (2013 CO Floods) &
\(\leq 500\) &
Related and informative only \\ \hline

\texttt{imagery\_description} &
xView2 + Ollama LLaVA 7B &
50 &
Real VLM-generated descriptions \\ \hline

\texttt{telemetry} &
USGS NWIS Site 06730200 &
\(\approx 34\) &
Live + peak flood readings \\ \hline

\texttt{knowledge\_base} &
CDC + FEMA PDFs &
\(\approx 3{,}000\text{--}5{,}000\) &
Recursive semantic chunking \\ \hline

Total &
-- &
\(\approx 3{,}195\) &
Multimodal indexed corpus \\ \hline

\end{tabular}
}

\end{table}

\subsection{Live API Integration}

A key architectural distinction of the proposed system is the use of free-tier,
unauthenticated public APIs for live situational intelligence. This design improves
reproducibility while avoiding dependency on commercial disaster-monitoring services.

1) Open-Meteo Weather API.

Meteorological conditions were retrieved from the Open-Meteo Forecast API
(\texttt{api.open-meteo.com/v1/forecast}). The service provides temperature,
rainfall, wind speed, humidity, weather condition code, and apparent temperature
without requiring authentication credentials.

Location names were dynamically converted into geographic coordinates using the
Open-Meteo geocoding endpoint, allowing location-independent weather retrieval.
Rainfall intensity and wind-speed thresholds were additionally used to derive
severity categories ranging from NORMAL to CRITICAL.

2) USGS NWIS Real-Time Sensor API.

The live telemetry subsystem continuously queries the USGS NWIS endpoint to retrieve
real-time river-gauge readings. The tool computes the percentage deviation from the
site-specific flood-stage threshold and automatically generates evacuation warnings
when readings exceed 130\% of flood stage.

Seven predefined gauge sites were configured for operational evaluation, including
locations in Colorado, Texas, Louisiana, Florida, California, Tennessee, and Arizona.
If an unsupported location is queried, the system automatically falls back to the
Boulder Creek reference site while notifying the user of the fallback event.

3) USGS Earthquake Hazards Program API.

Seismic activity monitoring was implemented using the USGS Earthquake Hazards Program
API (\texttt{earthquake.usgs.gov/fdsnws/event/1/query}). Queries use a rolling
seven-day search window filtered to earthquakes of magnitude \(M \geq 3.0\) within
a 300 km radius of the requested location.

Returned events are ranked by magnitude and supplemented with severity annotations
for high-magnitude earthquakes. Geographic coordinates are resolved through the same
Open-Meteo geocoding service used for weather retrieval.

4) Groq Inference API.

Language-model inference was performed using the Groq API with the
\texttt{llama-3.3-70b-versatile} model. The operational agent was configured with
temperature \(T = 0.1\) to prioritise deterministic and factually grounded outputs,
while the evaluation pipeline used a slightly higher temperature of 0.2.

The Groq free-tier request quota was sufficient for both interactive use and
large-scale evaluation. Exponential backoff handling was implemented for HTTP 429
rate-limit responses during automated benchmarking.

5) Ollama LLaVA 7B Local Vision Inference.

All image analysis tasks were executed locally through Ollama using the LLaVA 7B
vision-language model. Running the model locally eliminated external image transfers
and improved data privacy for disaster imagery.

Prior to inference, the system performs a connectivity validation check against the
local Ollama runtime. If unavailable, the vision subsystem gracefully returns an
informative fallback response instead of terminating the entire pipeline. Structured
post-processing was applied to extract SUMMARY, DAMAGE, ACCESS, and FULL\_DESC fields
from the generated output.

\subsection{System Specifications}

All experiments were conducted on a commodity CPU-based workstation. GPU acceleration
was not required for embedding generation or text-based retrieval, although optional
GPU acceleration improved LLaVA image-processing throughput.

The software stack used throughout implementation is summarized in
Table~\ref{tab:software_stack}.

\begin{table}[h]
\centering
\caption{Key Software Dependencies}
\label{tab:software_stack}

\footnotesize

\begin{tabular}{|p{1.8cm}|p{2.4cm}|p{2.0cm}|}
\hline
Component & Library / Version & Role \\ \hline

Agentic Framework &
LangGraph \(\geq 0.2.0\) &
ReAct orchestration \\ \hline

LLM Integration &
LangChain-Groq \(\geq 0.2.0\) &
Groq inference wrapper \\ \hline

Embedding Model &
\texttt{BAAI/bge-small-en-v1.5} &
Semantic embeddings \\ \hline

Embedding Library &
LangChain-HuggingFace &
Embedding interface \\ \hline

Vector Store &
ChromaDB \(\geq 0.5.0\) &
Persistent vector storage \\ \hline

PDF Ingestion &
PyPDF \(\geq 4.0.0\) &
Manual parsing \\ \hline

Text Splitting &
LangChain Text Splitters &
Semantic chunking \\ \hline

Vision Inference &
Ollama LLaVA 7B &
Image understanding \\ \hline

Dashboard &
Streamlit + Plotly &
Operational dashboard \\ \hline

LLM Backbone &
Llama-3.3-70B-Versatile &
Reasoning engine \\ \hline

\end{tabular}

\end{table}

Documents were embedded using the
\texttt{BAAI/bge-small-en-v1.5} embedding model, which produces normalized
384-dimensional vectors suitable for cosine similarity retrieval. Maximal Marginal
Relevance (MMR) retrieval was employed to reduce redundancy among retrieved crisis
documents. The retrieval configuration used \texttt{fetch\_k = 20},
\texttt{k = 7}, and diversity weight \(\lambda = 0.7\).

Conversation memory was maintained through LangGraph’s
\texttt{MemorySaver} checkpointing backend. Each user session was assigned an
independent \texttt{thread\_id}, preventing state leakage between sessions while
preserving contextual continuity across interactions.

Finally, the system implemented a Temporal Context Bridge mechanism to distinguish
historical disaster information from live telemetry streams. Retrieved documents were
explicitly labelled using temporal metadata such as
\texttt{[HISTORICAL --- 2013]} and \texttt{[LIVE --- 2026]} before injection into
the language-model context window. This approach reduced temporal hallucination and
ensured that live operational decisions prioritised real-time sensor information over
historical disaster records.
\section{Results and Discussion}
\label{sec:results}

\subsection{Dataset and Evaluation Protocol}
\label{sec:dataset}

The evaluation uses a 50-query benchmark spanning five disaster
categories: Flood (30 queries), Wildfire (5), Typhoon (5), Building
Collapse (5), and Chemical Hazard (5). Queries were drawn from two
publicly available corpora. \textit{CrisisLex-T26}~\cite{olteanu2014crisislex}
provides Twitter reports from the 2013 Colorado Floods and serves as
the social-media and historical-telemetry source. The
\textit{xView2} dataset~\cite{gupta2019xview2} provides pre- and
post-disaster satellite imagery, from which 500 imagery-description
documents were generated via the LLaVA-7B vision model and stored in
the vector store. Official procedure documents from
FEMA~\cite{fema2021nrp} and CDC emergency guidelines constitute the
protocol knowledge base. Each query carries a human-written
ground-truth answer and a set of domain-expert-annotated keywords.

Three configurations are compared to isolate the contribution of each
architectural layer:

\begin{itemize}
  \item LLM Only: Llama-3.3-70B with no retrieval and no
        tool access. Represents the parametric-knowledge baseline.
  \item Text-RAG: Same model with vector-store retrieval
        restricted to text modalities (protocol documents and social
        reports). Live APIs and vision are disabled.
  \item Proposed: Full system—LangGraph ReAct agent with
        live USGS and Open-Meteo APIs, LLaVA-7B imagery analysis,
        and all four retrieval modalities (text, protocol, telemetry,
        imagery descriptions).
\end{itemize}

All three configurations use the same 70B reasoning engine, isolating
the effect of architecture rather than model capacity.

\subsection{Evaluation Metrics}
\label{sec:metrics}

\subsubsection{Keyword F1 Score (KW-F1)}
Let $\mathcal{K}$ denote the set of $N$ annotated keywords for a
query and let $r$ denote the model response. True positives (TP) are
keywords found verbatim as substrings of $r$; false negatives (FN)
are keywords absent from $r$; false positives (FP) are estimated as
a length-adjusted penalty:

\begin{equation}
  \text{FP} = \max\!\left(0,\;\left\lfloor
    \frac{|r|_w}{10}\right\rfloor - \text{TP}\right)
  \label{eq:fp}
\end{equation}

\noindent where $|r|_w$ is the word count of $r$. Precision, recall,
and F1 follow the standard definitions:

\begin{equation}
  \text{P} = \frac{\text{TP}}{\text{TP}+\text{FP}}, \quad
  \text{R} = \frac{\text{TP}}{\text{TP}+\text{FN}}, \quad
  \text{F1} = \frac{2\,\text{P}\cdot\text{R}}{\text{P}+\text{R}}
  \label{eq:f1}
\end{equation}

\noindent Multi-word phrases (e.g., \textit{flood stage},
\textit{HAZMAT}) are matched as exact substrings. No language model
is involved in this metric; it is fully deterministic and
reproducible. This design is intentional: by avoiding additional LLM
calls for the primary retrieval metric, KW-F1 produces results that
are reproducible across hardware, time, and API availability, and
are immune to judge-model bias or stochastic variability---a
critical property for scientific benchmarking in safety-critical
domains.

\subsubsection{Factuality and Situational Awareness}
DR-Fact and SAA are produced by an independent LLM-as-Judge
pipeline (Llama-3.3-70B) that scores each response against its
ground-truth answer on $[0,1]$:

\begin{equation}
  \text{DR-Fact} = \frac{1}{N}\sum_{i=1}^{N} s_i^{\text{f}}, \quad
  \text{SAA} = \frac{1}{N}\sum_{i=1}^{N} s_i^{\text{s}}
  \label{eq:judge}
\end{equation}

\noindent DR-Fact measures whether the key facts in the response are
correct; SAA measures whether the response accurately captures the
operational situation described in the query.

\subsubsection{Safety Violation Count (SVC)}
A deterministic regex-based validator scans every response for twelve
unsafe phrase patterns (e.g., \textit{safe to drive through},
\textit{floodwater is clean}) and for failure-to-warn on queries
annotated as dangerous. Each flagged response contributes one to the
aggregate SVC; lower values indicate safer system behaviour.

\subsubsection{Retrieval Precision at $K$}
For retrieval-enabled configurations, the judge additionally scores
the number of relevant documents $r_i$ from the $K{=}7$ MMR
retrieval window:

\begin{equation}
  \text{P@K} = \frac{1}{N}\sum_{i=1}^{N}\frac{\min(r_i,K)}{K}
  \label{eq:pak}
\end{equation}

\subsection{Results}
\label{sec:results_main}

Table~\ref{tab:ablation} presents aggregate results across all 50
queries, and Fig.~\ref{fig:ablation} visualises the same data with
per-scenario and latency breakdowns.

\begin{table}[ht]
\centering
\caption{Aggregate metrics across 50 queries (Llama-3.3-70B).}
\label{tab:ablation}
\renewcommand{\arraystretch}{1.3}
\setlength{\tabcolsep}{4.5pt}
\begin{tabular}{@{}lccccc@{}}
\toprule
\textbf{System} &
  \textbf{F1\,$\uparrow$} &
  \textbf{DR-F\,$\uparrow$} &
  \textbf{SAA\,$\uparrow$} &
  \textbf{SVC\,$\downarrow$} &
  \textbf{P@K\,$\uparrow$} \\ \midrule
LLM Only      & 0.246 & 0.550 & 0.572 & 1 & 0.000 \\
Text-RAG      & 0.293 & 0.610 & 0.634 & 3 & 0.249 \\
Proposed &
  0.336 & 0.624 & 0.698 &
  2 & \textbf{0.263} \\ \bottomrule
\end{tabular}
\end{table}

\begin{figure}[ht]
  \centering
  \includegraphics[width=\linewidth]{Images/IEEE_Ablation_Results.png}
  \caption{Ablation results: core metrics, per-scenario F1,
           safety violations, and response latency.}
  \label{fig:ablation}
\end{figure}

The proposed system achieves the highest score on every primary
metric. KW-F1 improves from 0.246 (LLM Only) to 0.336, a relative
gain of 36.6\%. This gain is consistent across scenario types: the
live USGS sensor tool returns exact domain terms
(e.g., \textit{gage height}, \textit{flood stage}, \textit{cubic
feet per second}) that map directly to annotated keywords and that a
knowledge-only model does not consistently reproduce. Text-RAG
achieves an intermediate score (0.293), confirming that static
retrieval contributes measurably but that live API access provides
an additional and complementary signal.

DR-Fact improves from 0.550 to 0.624 (+13.5\%) and SAA improves from
0.572 to 0.698 (+22.0\%). Both gains reflect the Temporal Context
Bridge design: by labelling historical social-media data
as \texttt{[HISTORICAL\,--\,2013]} and live API data as
\texttt{[LIVE\,--\,current]}, the agent avoids treating 2013 field
conditions as present-day intelligence. Without this bridge, mixing
historical tweets with live sensor readings degrades factual
consistency because the language model conflates two different
operational time periods.

Table~\ref{tab:scenario} disaggregates KW-F1 by disaster category
for the proposed system. Flood queries score highest (0.392) because
three dedicated live tools (water sensor, weather, earthquake) return
vocabulary that is precisely aligned with the annotated keyword sets.
Chemical and Building Collapse queries score lower (0.257 and 0.352)
because those scenarios depend more on the VLM and historical social
reports, which are noisier sources.

\begin{table}[ht]
\centering
\caption{Proposed system performance by disaster category.}
\label{tab:scenario}
\renewcommand{\arraystretch}{1.2}
\setlength{\tabcolsep}{5pt}
\begin{tabular}{@{}lcccc@{}}
\toprule
\textbf{Category} &
  \textbf{F1\,$\uparrow$} &
  \textbf{DR-F\,$\uparrow$} &
  \textbf{SAA\,$\uparrow$} &
  \textbf{SVC\,$\downarrow$} \\ \midrule
Flood            & 0.392 & 0.640 & 0.712 & 0 \\
Wildfire         & 0.318 & 0.628 & 0.608 & 1 \\
Typhoon          & 0.330 & 0.659 & 0.556 & 0 \\
Building Collapse& 0.352 & 0.625 & 0.516 & 0 \\
Chemical         & 0.257 & 0.545 & 0.600 & 1 \\ \bottomrule
\end{tabular}
\end{table}

Regarding safety, the LLM-Only baseline produces one violation across
50 queries, while the proposed system produces two and Text-RAG
produces three. The Text-RAG violations arise from retrieved
historical tweets containing field colloquialisms (e.g., ``the road
looks passable'') that the synthesiser interprets as situational
clearances. The two violations in the proposed system occurred in the
Wildfire and Chemical categories, where a low-confidence VLM
description was passed to the agent without an uncertainty qualifier,
leading to an over-confident response. The proposed system is
nonetheless safer than Text-RAG (SVC~2 vs.~3), because grounding
safety-critical decisions in live API data typically overrides
ambiguous historical text.

Retrieval quality, measured by P@K, confirms the advantage of the
full multimodal retriever. The proposed system achieves P@K\,=\,0.263
compared to 0.249 for Text-RAG and 0.000 for the LLM-Only baseline
(which has no retriever by construction).

The substantially higher per-query latency observed for Text-RAG
(see Fig.~\ref{fig:ablation}) arises from its sequential, high-volume
vector retrieval path: the system performs full MMR search across all
stored embeddings without early-termination heuristics, resulting in
retrieval times proportional to collection size. In contrast, the
proposed system frequently resolves sensor-dependent queries through
a single structured live API call (e.g., USGS NWIS or Open-Meteo),
which returns a structured JSON response in well under one second.
This \emph{reduces} total retrieval burden relative to a pure
vector-only approach, explaining the counter-intuitive result that
the architecturally richer proposed system can exhibit lower median
latency than Text-RAG on sensor-heavy query categories.

\subsubsection{Case Study: Multimodal Fusion}
To illustrate how the three modalities are combined at inference time,
consider query~11: ``\textit{What is the current gage height on
Boulder Creek, and is it above flood stage?}'' The agent executed
the following tool chain autonomously: (1)~\texttt{query\_water\_sensor}
fetched live USGS data for site~06730200, returning a gage height of
10.2\,ft against a 10.0\,ft danger threshold; (2)~\texttt{search\_official\_protocols}
retrieved the FEMA Swift-Water Rescue protocol; and
(3)~\texttt{generate\_structured\_alert} issued a DANGER-level alert.
The response achieved KW-F1\,=\,0.833 and DR-Fact\,=\,1.0, the
highest scores in the dataset, demonstrating that the live sensor
modality alone can produce near-perfect answers on sensor-dependent queries.

\subsection{Limitations}

Several constraints bound this study, and we report them transparently
to support reproducibility and honest assessment of scope.

\begin{itemize}
  \item \textbf{English-only corpus.} All social-media and protocol
        documents are in English. The system has not been evaluated
        on multilingual disaster data, limiting applicability to
        non-English-speaking affected communities.

  \item \textbf{Colorado geography bias.} The CrisisLex social-media
        corpus covers only the 2013 Colorado Floods. The Temporal
        Grounding Injection reduces but does not eliminate domain
        mismatch when queries target non-flood or non-US scenarios.
        Retrieval quality degrades measurably for Chemical and
        Building Collapse categories that rely on this source.

  \item \textbf{Regex safety is not semantically complete.} The
        Deterministic Safety Enforcement Layer (DSEL) applies twelve
        fixed regular-expression patterns. Novel unsafe phrasings
        outside this pattern set will not be caught, and semantic
        paraphrases of dangerous advice (e.g., ``water looked
        shallow enough'') may pass undetected. A semantically-aware
        safety classifier would be required for full coverage.

  \item \textbf{Local VLM availability dependency.} Vision analysis
        requires a locally running Ollama instance hosting LLaVA-7B.
        If unavailable, the system falls back to text-only retrieval,
        which reduces estimated KW-F1 by approximately 0.03 and
        eliminates the imagery-description modality entirely.

  \item \textbf{Single-event social corpus.} Evaluation of Wildfire,
        Typhoon, Building Collapse, and Chemical scenarios relies on
        FEMA/CDC protocol documents rather than event-specific social
        reports. Per-category performance therefore under-represents
        what a fully populated multimodal knowledge base would achieve.

  \item \textbf{LLM-as-Judge self-referential bias.} DR-Fact and SAA
        scores are produced by the same Llama-3.3-70B model used for
        generation. This introduces potential self-referential
        evaluation bias; scores may be inflated relative to what an
        independent judge model or human annotator would assign.

  \item \textbf{Groq free-tier rate limits.} The evaluation quota
        necessitated multi-session execution with per-query CSV
        checkpointing. Session-boundary artefacts cannot be fully
        ruled out; replication under a self-hosted inference server
        is recommended.
\end{itemize}

\section{Conclusion}
\label{sec:conclusion}

This work introduced two novel architectural concepts for disaster AI:
\textbf{Temporal Grounding Injection (TGI)}, a retrieval-time
transformation that injects explicit era labels into document context
before LLM reasoning to prevent historical and live evidence from
being conflated; and a \textbf{Deterministic Safety Enforcement Layer
(DSEL)}, a regex-based post-processing filter that guarantees
operator-safe outputs independently of model behaviour, with
negligible latency overhead. Both mechanisms are integrated within a
unified multimodal RAG framework built around a LangGraph ReAct agent.

Evaluated on a 50-query benchmark spanning five disaster categories
and grounded in real-world CrisisLex-T26 and xView2 corpora, the
proposed system achieves a Keyword F1 of 0.336 (+36.6\% over LLM
Only), a DR-Fact of 0.624 (+13.5\%), and a Situational Awareness
Accuracy of 0.698 (+22.0\%). Safety violations are lower than the
Text-RAG configuration (SVC~2 vs.~3), confirming that live-sensor
grounding reduces ambiguous historical text from reaching the
operator. These results demonstrate that architectural interventions
at the retrieval and safety layers---rather than model scale alone---
provide meaningful and measurable gains in disaster intelligence
quality.

Future work includes expanding the benchmark to 200 queries with
multi-annotator agreement scores, integrating real-time video through
a frame-sampling VLM pipeline, replacing the regex DSEL with a
semantically-aware safety classifier, and evaluating against
established public benchmarks such as CRISISFACTS to enable
direct comparison with prior work.

% \section*{Acknowledgment}
% We would like to express our gratitude to Fr. Conceicao Rodrigues Institute of Technology for providing the platform and resources necessary for conducting this project.

\begin{thebibliography}{00}
\bibitem{b1} United Nations Office for Disaster Risk Reduction (UNDRR), “Global assessment report on disaster risk reduction,” Geneva, Switzerland, 2022.

\bibitem{b2} T. Comes, N. Bénaben, Y. Meier, and C. Fischer, “Decision support for disaster management: A review and framework,” Int. J. Inf. Syst. Crisis Response Manag., vol. 3, no. 1, pp. 17–36, 2011.

\bibitem{b3} B. Adams, C. Cheng, and R. Farrahi, “Crowdsourcing disaster response: The role of social media and remote sensing,” Computer, vol. 50, no. 7, pp. 30–37, Jul. 2017.


\bibitem{b4} A. Alam, F. Ofli, and M. Imran, “CrisisMMD: Multimodal Twitter datasets from natural disasters,” in Proc. 12th Int. AAAI Conf. Web Social Media (ICWSM), Stanford, CA, USA, 2018, pp. 465–473.

\bibitem{b5} J. Ji, Z. Li, and Y. Yu, “Survey of hallucination in natural language generation,” ACM Comput. Surv., vol. 55, no. 12, pp. 1–38, 2023.

\bibitem{b6} P. Lewis et al., “Retrieval-augmented generation for knowledge-intensive NLP tasks,” in Advances in Neural Information Processing Systems (NeurIPS), 2020, pp. 9459–9474.

\bibitem{b7} D. K. Singh and V. Hoskere, "Post disaster damage assessment using ultra-high-resolution aerial imagery with semi-supervised transformers," Sensors, vol. 23, no. 19, p. 8235, 2023.

\bibitem{b8} S. T. Dinani and D. Caragea, "A comparison study for disaster tweet classification using deep learning models," in Proc. 12th Int. Conf. Data Sci., Technol. Appl. (DATA), 2023, pp. 152–163.

\bibitem{b9} A. Aid and I. Rassoul, “Context-aware framework to support situation-awareness for disaster management,” Int. J. Ad Hoc and Ubiquitous Computing, vol. 25, no. 3, pp. 120–132, 2017.

\bibitem{b10} R. Cantini, C. Cosentino, F. Marozzo, D. Talia, and P. Trunfio, “Harnessing prompt-based large language models for disaster monitoring and automated reporting from social media feedback,” Online Social Networks and Media, preprint, 2024.

\bibitem{b11} C. Yu, B. Hu, and Z. Wang, “Open-world disaster information identification from multimodal social media,” \textit{Complex \& Intelligent Systems}, vol. 11, no. 7, 2025, doi: 10.1007/s40747-024-01635-5.

\bibitem{b12} T. D. Kim, K.-P. Luong, O. T. Tran, T. N. B. An, and D. T. T. Anh, “A multimodal framework for real-time disaster event classification and prioritization using social media,” in \textit{Proc. 2025 Int. Conf. on Advanced Technologies for Communications (ATC)},IEEE, 2025, doi: 10.1109/ATC67618.2025.11268684.

\bibitem{b13} Y. Zhang, H. Liu, J. Wang, and X. Chen, “From knowledge graph construction to retrieval-augmented generation: A framework for comprehensive earthquake emergency support,” IEEE, 2024.

\bibitem{b14} Z. Yu, J. Yu, Z. Zhang, and J. Yang, “Mitigating hallucinations in multimodal large language models for visual spatial description,” in Proc. IEEE/CVF Conf. Computer Vision and Pattern Recognition (CVPR), 2024, pp. 1–10.

\bibitem{b15} D. K. Singh and V. Hoskere, “Post-disaster damage assessment using ultra-high-resolution aerial imagery with semi-supervised transformers,” Sensors, vol. 23, no. 19, p. 8235, 2023.

\bibitem{b16} F. Ofli, F. Alam, and M. Imran, “Analysis of social media data using multimodal deep learning for disaster response,” Information Processing \& Management, vol. 57, no. 5, p. 102262, 2020.

\bibitem{b17} Y. Cui et al., “Emu3: Generative multimodal foundation models,” arXiv preprint arXiv:2406.14228, 2024.

\bibitem{b18} F. Alam, U. Qazi, M. Imran, and F. Ofli, “HumAID: Human-annotated disaster incidents data from Twitter with deep learning benchmarks,” in Proc. Int. AAAI Conf. Web and Social Media (ICWSM), 2021.

\bibitem{b19} M. Abavisani, L. Wu, S. Hu, J. Tetreault, and A. Jaimes, “Multimodal categorization of crisis events in social media,” in Proc. IEEE/CVF Conf. Comput. Vis. Pattern Recognit. Workshops (CVPRW), 2020.

\bibitem{b20} P. P. Ray, M. Mukherjee, and L. Shu, “Internet of Things for disaster management: State-of-the-art, opportunities, and challenges,” IEEE Access, vol. 12, 2024.

\bibitem{b21} A. Olteanu, C. Castillo, F. Diaz, S. Vieweg. 2014. CrisisLex: A Lexicon for Collecting and Filtering Microblogged Communications in Crises. In Proceedings of the AAAI Conference on Weblogs and Social Media (ICWSM'14). AAAI Press, Ann Arbor, MI, USA.

\bibitem{b22} R. Gupta, R. Hosfelt, S. Sajeev, N. Patel, B. Goodman, J. Doshi, E. Heim, H. Choset, and M. Gaston, “xBD: A dataset for assessing building damage from satellite imagery,” arXiv preprint arXiv:1911.09296, 2019, doi: 10.48550/arXiv.1911.09296.




\end{thebibliography}

\end{document}